% Generated by src/11_attention.py. Requires \usepackage{booktabs}.
\begin{table}[htbp]
\centering
\small
\caption{Investor attention on the reaction day: Wikipedia pageviews}
\label{tab:attention}
\begin{tabular}{lccc}
\toprule
 & Treated & Control & Difference \\
\midrule
Abnormal views, mean (log points) & 0.214 & 0.271 & $-$0.057 \\
 & & & $($-$9.32)$ \\
Abnormal views, median & 0.166 & 0.210 & \\
Implied change in views vs.\ baseline (\%) & 23.8 & 31.1 & \\
Baseline daily views, median & 110 & 120 & \\
\midrule
Matched pairs: mean difference & & & $-$0.062 \\
 & & & $($-$8.70)$ \\
Regression on MegaDay (clustered by day) & & & $-$0.057 \\
 & & & $($-$4.41)$ \\
\midrule
\multicolumn{4}{l}{\emph{MegaDay coefficient with controls}} \\
All mega-cap days & & & $-$0.032 \\
 & & & $($-$4.02)$ \\
Nvidia days only & & & $-$0.013 \\
 & & & $($-$0.71)$ \\
Other five mega-caps & & & $-$0.032 \\
 & & & $($-$3.95)$ \\
Same-event matches only & & & $-$0.020 \\
 & & & $($-$2.16)$ \\
\midrule
Announcements with a Wikipedia page & 7,439 & 7,399 & \\
Matched pairs with both pages & & & 4,594 \\
Share of sample with a page & 54\% & 54\% & \\
\bottomrule
\end{tabular}
\par\medskip
\parbox{0.95\linewidth}{\footnotesize Abnormal views $=\log(1+v_0)-\log(1+\tilde v)$, where $v_0$ is English Wikipedia pageviews (users, all access) on the reaction day and $\tilde v$ is the median daily views over the 56 days ending 3 days before it. Wikipedia pages are matched to firms through the SEC CIK recorded in Wikidata (P5531). 1:1 matched sample. $t$-statistics in parentheses: Welch test for the difference in means, paired test for matched pairs, and standard errors clustered by reaction day for the regressions. Controls: day-of-week, calendar-quarter and FF12 industry fixed effects, log baseline views, log market cap, log(1 + analysts) and the same-day announcement count. Nvidia and other-mega-cap rows compare the treated firms on those days with all controls. Firms with a Wikipedia page are larger and better known than the average sample firm. Crowding out predicts lower attention to treated firms; crowding in predicts higher.}
\end{table}
